% !TeX root = ../navier-stokes-manuscript.tex

We prove global smoothness for the unforced three-dimensional incompressible Navier--Stokes equations on \(\R^3\) for every viscosity \(\nu>0\) and every divergence-free Schwartz initial velocity \(u_0\).
The corresponding velocity and pressure remain smooth for all \(t\geq 0\).

The proof rules out a finite singular endpoint directly.
Let \([0,T_*)\) be the maximal classical lifespan and suppose \(T_*<\infty\).
Repeated time differentiation of the Navier--Stokes equations, with pressure recovered at each order by Riesz transforms, generates one pressure-complete mixed jet of the same solution.
For each fixed finite temporal depth \(N\) and spatial Sobolev height \(M\), the central estimate controls the corresponding rectangle of that jet uniformly over every cutoff \(T<T_*\).
The adjacent-row trace estimate then gives, for every finite \(r\) and \(m\),
\[
\sup_{0\leq t<T_*}
\|\partial_t^r u(t)\|_{H^m(\R^3)}
<\infty.
\]

These finite rectangles agree under restriction.
Their compatible intersection preserves the whole-terminal control at every finite order, and the same trace mechanism gives a strong terminal limit in every Sobolev space:
\[
u(t)\longrightarrow u_*
\quad\text{in }H^m(\R^3)
\quad\text{as }t\uparrow T_*
\]
for every finite \(m\), with one terminal velocity
\[
u_*\in\bigcap_{m=0}^{\infty}H^m(\R^3).
\]

The differentiated Navier--Stokes equations and the pressure recovery determine the complete velocity-pressure jet at \(T_*\).
Thus the alleged singular endpoint is a smooth time slice of the same solution.
Smooth local existence restarts the system from \(u_*\), and uniqueness identifies that restart with the original solution.
The same solution then continues beyond \(T_*\), contradicting maximality.
Hence \(T_*=\infty\).
